\chapter{THỰC HIỆN CÀI ĐẶT}

\section{Cấu trúc tổ chức mã nguồn}

\subsection{Repository Structure}
Mã nguồn dự án trên GitHub được phân chia thành các module/thư mục rõ ràng:

\begin{verbatim}
Si_Sap_InteractHub-main/
├── /backend                  # ASP.NET Core API
│   ├── Controllers/          # API endpoints
│   ├── Services/             # Business logic
│   ├── Models/
│   │   ├── Entities/         # Database entities
│   │   └── DTOs/             # Data transfer objects
│   ├── Data/                 # DbContext
│   ├── Migrations/           # EF Core migrations
│   ├── Program.cs            # Startup configuration
│   ├── appsettings.json      # Configuration
│   └── backend.csproj        # Project file
│
├── /frontend                 # React TypeScript SPA
│   ├── src/
│   │   ├── app/
│   │   │   ├── App.tsx       # Main component
│   │   │   ├── routes.tsx    # Route definitions
│   │   │   └── components/   # Reusable components
│   │   ├── pages/            # Page components
│   │   ├── services/         # API services
│   │   ├── hooks/            # Custom React hooks
│   │   ├── contexts/         # Context API
│   │   ├── types/            # TypeScript types
│   │   └── styles/           # Global styles
│   ├── vite.config.ts        # Vite bundler config
│   ├── tailwind.config.mjs   # Tailwind CSS config
│   ├── package.json          # Dependencies
│   └── index.html
│
├── /database                 # Database scripts
│   └── InteractHub_full.sql  # Complete database schema
│
├── /BaoCao                   # Report (LaTeX)
│   ├── chapters/             # Report chapters
│   ├── images/               # Report images
│   ├── main.tex              # Main report file
│   └── main.pdf              # Generated PDF
│
├── /guidelines               # Documentation
│   └── Guidelines.md         # Project guidelines
│
├── README.md                 # Project readme
├── backend.sln               # Visual Studio solution
└── InteracHub.slnx           # Alternative solution format
\end{verbatim}

\section{Phát triển Backend}

\subsection{Database Setup}

\subsubsection{Entity Framework Core Configuration}
\begin{verbatim}
// Program.cs
builder.Services.AddDbContext<AppDbContext>(options =>
    options.UseSqlServer(
        builder.Configuration.GetConnectionString(
            "DefaultConnection")));
\end{verbatim}

\subsubsection{Migrations}
Dự án sử dụng Code-First migrations để quản lý schema:

\begin{itemize}
    \item \textbf{InitialCreate:} Tạo schema ban đầu với 10 entities
    \item \textbf{AddEntityConstraintsAndSeedHashtags:} Thêm constraints và seed data hashtags phổ biến
    \item \textbf{AddQueryIndexes:} Tối ưu query performance bằng indexes
\end{itemize}

Command tạo migrations:
\begin{verbatim}
Add-Migration [MigrationName]
Update-Database
\end{verbatim}

\subsection{API Controllers}

\subsubsection{AuthController}
Xử lý xác thực:

\begin{verbatim}
POST   /api/auth/register         - Đăng ký user mới
POST   /api/auth/login            - Login & nhận JWT token
POST   /api/auth/logout           - Logout
POST   /api/auth/refresh-token    - Refresh JWT token
\end{verbatim}

Features:
\begin{itemize}
    \item Password validation (minimum 6 characters)
    \item Email uniqueness check
    \item JWT token generation
    \item Secure password hashing via Identity
\end{itemize}

\subsubsection{PostsController}
Quản lý bài viết:

\begin{verbatim}
GET    /api/posts                 - Lấy danh sách posts (paging)
GET    /api/posts/{id}            - Chi tiết post
POST   /api/posts                 - Tạo post mới
PUT    /api/posts/{id}            - Cập nhật post
DELETE /api/posts/{id}            - Xóa post
GET    /api/posts/{id}/comments   - Danh sách comments
GET    /api/posts/{id}/likes      - Danh sách likes
\end{verbatim}

Features:
\begin{itemize}
    \item Pagination (skip, take)
    \item Include related data (User, Comments, Likes)
    \item Image upload to Azure Blob Storage
    \item Hashtag support
    \item Only owner can edit/delete own posts
\end{itemize}

\subsubsection{UsersController}
Quản lý user profiles:

\begin{verbatim}
GET    /api/users/{id}            - User profile
PUT    /api/users/{id}            - Cập nhật profile
GET    /api/users/{id}/posts      - Posts của user
GET    /api/users/{id}/friends    - Danh sách bạn bè
\end{verbatim}

\subsubsection{FriendsController}
Quản lý quan hệ kết bạn:

\begin{verbatim}
POST   /api/friends/request          - Gửi friend request
PUT    /api/friends/{id}/accept      - Accept request
PUT    /api/friends/{id}/reject      - Reject request
DELETE /api/friends/{id}             - Remove friend
GET    /api/friends                  - Danh sách bạn bè
\end{verbatim}

\subsubsection{Các Controllers khác}
\begin{itemize}
    \item \textbf{StoriesController:} Quản lý stories (24 giờ expiry)
    \item \textbf{NotificationsController:} Quản lý thông báo real-time
\end{itemize}

\subsection{Service Layer}

Các services chứa business logic:

\begin{verbatim}
IPostsService / PostsService
├── CreatePost()
├── UpdatePost()
├── DeletePost()
├── GetPostById()
├── GetUserPosts()
└── GetPostWithComments()

IUsersService / UsersService
├── GetUserProfile()
├── UpdateProfile()
├── GetUserById()
└── SearchUsers()

IFriendsService / FriendsService
├── SendFriendRequest()
├── AcceptRequest()
├── RejectRequest()
├── GetFriendsList()
└── RemoveFriend()

INotificationsService / NotificationsService
├── CreateNotification()
├── GetUserNotifications()
├── MarkAsRead()
└── DeleteNotification()
\end{verbatim}

Features:
\begin{itemize}
    \item Validation logic
    \item Authorization checks (verify ownership)
    \item Error handling
    \item Async/await for performance
\end{itemize}

\subsection{DTOs (Data Transfer Objects)}

DTOs được sử dụng để:
\begin{itemize}
    \item Tách biệt entities từ API response
    \item Validate input từ client
    \item Control fields được expose
\end{itemize}

\begin{verbatim}
AuthDtos.cs
├── RegisterRequest
├── LoginRequest
├── LoginResponse
└── JwtTokenResponse

PostDtos.cs
├── CreatePostRequest
├── UpdatePostRequest
├── PostResponse
├── PostListResponse
└── PostDetailResponse

SocialDtos.cs
├── CommentRequest
├── CommentResponse
├── LikeResponse
├── FriendshipRequest
└── FriendshipResponse
\end{verbatim}

\subsection{Authentication \& Authorization}

\subsubsection{JWT Setup}
\begin{verbatim}
// Program.cs
var jwtKey = builder.Configuration["Jwt:Key"];
var jwtIssuer = builder.Configuration["Jwt:Issuer"];
var jwtAudience = builder.Configuration["Jwt:Audience"];

builder.Services.AddAuthentication(JwtBearerDefaults.AuthenticationScheme)
    .AddJwtBearer(options => {
        options.TokenValidationParameters = new TokenValidationParameters {
            ValidateIssuer = true,
            ValidateAudience = true,
            ValidateIssuerSigningKey = true,
            ValidIssuer = jwtIssuer,
            ValidAudience = jwtAudience,
            IssuerSigningKey = new SymmetricSecurityKey(
                Encoding.UTF8.GetBytes(jwtKey))
        };
    });
\end{verbatim}

\subsubsection{Role-based Authorization}
\begin{verbatim}
// Program.cs
builder.Services.AddAuthorization(options => {
    options.AddPolicy("AdminOnly", policy => 
        policy.RequireRole("Admin"));
    
    options.AddPolicy("UserOrAdmin", policy => 
        policy.RequireRole("User", "Admin"));
});

// Controller usage
[Authorize(Policy = "AdminOnly")]
public async Task<IActionResult> DeletePost(int id) { ... }
\end{verbatim}

\section{Phát triển Frontend}

\subsection{Project Structure}

\begin{verbatim}
src/
├── app/
│   ├── App.tsx               # Main component wrapper
│   ├── routes.tsx            # Route configuration
│   └── components/
│       ├── Navbar.tsx
│       ├── Sidebar.tsx
│       ├── PostCard.tsx
│       ├── CommentSection.tsx
│       └── AuthForm.tsx
│
├── pages/
│   ├── HomePage.tsx
│   ├── ProfilePage.tsx
│   ├── FriendsPage.tsx
│   ├── NotificationsPage.tsx
│   └── LoginPage.tsx
│
├── services/
│   ├── apiClient.ts          # Axios instance with JWT
│   ├── authService.ts
│   ├── postService.ts
│   ├── userService.ts
│   └── notificationService.ts
│
├── hooks/
│   ├── useAuth.ts            # Auth context hook
│   ├── usePosts.ts           # Post fetching hook
│   └── useFriends.ts         # Friendship hook
│
├── contexts/
│   ├── AuthContext.tsx       # Auth state & functions
│   ├── NotificationContext.tsx
│   └── UserContext.tsx
│
├── types/
│   ├── User.ts
│   ├── Post.ts
│   ├── Comment.ts
│   ├── Friendship.ts
│   └── Notification.ts
│
└── styles/
    ├── index.css
    ├── tailwind.css
    └── theme.css
\end{verbatim}

\subsection{Component Architecture}

\subsubsection{React Hooks Usage}
\begin{verbatim}
// Custom hook example: useAuth.ts
export const useAuth = () => {
    const context = useContext(AuthContext);
    
    const login = async (email, password) => {
        const response = await apiClient.post('/auth/login', 
            { email, password });
        localStorage.setItem('token', response.data.token);
        return response.data;
    };
    
    const logout = () => {
        localStorage.removeItem('token');
        context.setUser(null);
    };
    
    return { ...context, login, logout };
};
\end{verbatim}

\subsubsection{Protected Routes}
\begin{verbatim}
// routes.tsx
const ProtectedRoute: React.FC<{element: JSX.Element}> = 
    ({ element }) => {
    const { user } = useAuth();
    return user ? element : <Navigate to="/login" />;
};

export const routes = [
    { path: '/', element: <HomePage /> },
    { path: '/login', element: <LoginPage /> },
    { path: '/profile/:id', 
      element: <ProtectedRoute element={<ProfilePage />} /> },
    { path: '/friends', 
      element: <ProtectedRoute element={<FriendsPage />} /> },
];
\end{verbatim}

\subsection{State Management}

\subsubsection{React Context API}
\begin{verbatim}
// contexts/AuthContext.tsx
export const AuthContext = createContext<AuthContextType | null>(null);

export const AuthProvider: React.FC<{children: ReactNode}> = 
    ({ children }) => {
    const [user, setUser] = useState<User | null>(null);
    const [token, setToken] = useState<string | null>(
        localStorage.getItem('token'));
    
    useEffect(() => {
        if (token) {
            apiClient.defaults.headers.common['Authorization'] = 
                `Bearer ${token}`;
            // Verify token & load user
        }
    }, [token]);
    
    return (
        <AuthContext.Provider value={{ user, token, setUser, setToken }}>
            {children}
        </AuthContext.Provider>
    );
};
\end{verbatim}

\subsection{API Communication}

\subsubsection{API Client Setup}
\begin{verbatim}
// services/apiClient.ts
import axios from 'axios';

export const apiClient = axios.create({
    baseURL: process.env.REACT_APP_API_URL || 
             'https://api.interacthub.com/api',
    timeout: 10000,
});

// Request interceptor (add JWT token)
apiClient.interceptors.request.use((config) => {
    const token = localStorage.getItem('token');
    if (token) {
        config.headers.Authorization = `Bearer ${token}`;
    }
    return config;
});

// Response interceptor (handle 401)
apiClient.interceptors.response.use(
    (response) => response,
    (error) => {
        if (error.response?.status === 401) {
            localStorage.removeItem('token');
            window.location.href = '/login';
        }
        return Promise.reject(error);
    }
);
\end{verbatim}

\subsection{Styling}

\subsubsection{Tailwind CSS}
Dự án sử dụng Tailwind CSS utility classes:

\begin{verbatim}
// Example: PostCard component
<div className="bg-white rounded-lg shadow-md p-4 
                hover:shadow-lg transition-shadow">
    <div className="flex items-center mb-4">
        <img src={post.user.avatar} 
             className="w-10 h-10 rounded-full mr-3" />
        <div>
            <h3 className="font-bold">{post.user.name}</h3>
            <p className="text-sm text-gray-500">
                {formatDate(post.createdAt)}
            </p>
        </div>
    </div>
    
    <p className="text-gray-800 mb-3">{post.content}</p>
    
    {post.imageUrl && (
        <img src={post.imageUrl} 
             className="w-full rounded-lg mb-3 
                       max-h-96 object-cover" />
    )}
    
    <div className="flex justify-between 
                    text-gray-600 text-sm border-t pt-3">
        <button className="hover:text-blue-500">
            ❤️ {post.likes.length} Likes
        </button>
        <button className="hover:text-blue-500">
            💬 {post.comments.length} Comments
        </button>
    </div>
</div>
\end{verbatim}

\section{Configuration}

\subsection{appsettings.json (Backend)}
\begin{verbatim}
{
  "ConnectionStrings": {
    "DefaultConnection": "Server=[server];Database=InteractHub;
                         User Id=[user];Password=[password];"
  },
  "Jwt": {
    "Key": "[your-secret-key-min-32-chars]",
    "Issuer": "InteractHub.API",
    "Audience": "InteractHub.Client",
    "ExpirationMinutes": 60
  },
  "AzureStorage": {
    "ConnectionString": "[azure-storage-connection-string]",
    "ContainerName": "images"
  },
  "Logging": {
    "LogLevel": {
      "Default": "Information"
    }
  }
}
\end{verbatim}

\subsection{.env.local (Frontend)}
\begin{verbatim}
REACT_APP_API_URL=https://api.interacthub.com/api
REACT_APP_SIGNALR_HUB_URL=https://api.interacthub.com/notificationHub
REACT_APP_ENVIRONMENT=production
\end{verbatim}

\section{Build \& Deployment}

\subsection{Backend Build}
\begin{verbatim}
# Development
dotnet run

# Production build
dotnet publish -c Release -o ./publish

# Create docker image
docker build -t interacthub-api:latest .
docker push [registry]/interacthub-api:latest
\end{verbatim}

\subsection{Frontend Build}
\begin{verbatim}
# Install dependencies
npm install

# Development server
npm run dev

# Production build
npm run build    # Creates dist/ folder
npm run preview  # Preview production build
\end{verbatim}
